\documentclass[a4paper,11pt,twocolumn]{ctexart}

% === 页面设置 ===
\usepackage[top=2.5cm,bottom=2.5cm,left=2cm,right=2cm]{geometry}
\usepackage{setspace}
\setstretch{1.2}

% === 字体 ===
\ctexset{
    section/format = \zihao{4}\heiti,
    subsection/format = \zihao{-4}\heiti,
    subsubsection/format = \zihao{-4}\kaishu,
}

% === 常用宏包 ===
\usepackage{amsmath,amssymb,amsthm}
\usepackage{graphicx}
\usepackage{booktabs}
\usepackage{multirow}
\usepackage{float}
\usepackage{caption}
\usepackage{subcaption}
\captionsetup{labelsep=space}  % "图 1 xxx" 而非 "图 1: xxx"
\usepackage[hidelinks]{hyperref}
\usepackage{cleveref}
\usepackage{algorithm2e}
\usepackage{enumitem}
\usepackage{abstract}

% === TikZ（仅用于架构图/流程图）===
\usepackage{tikz}
\usetikzlibrary{arrows.meta,positioning,shapes.geometric,fit,calc}
% 注意：pgfplots 仅作为备用，数据图表应使用 \includegraphics 引用 matplotlib 生成的 PDF
\usepackage{pgfplots}
\pgfplotsset{compat=1.18}

% === 引用格式 ===
\usepackage[super,sort&compress]{gbt7714}

% === 定理环境 ===
\newtheorem{definition}{定义}[section]
\newtheorem{theorem}{定理}[section]
\newtheorem{lemma}{引理}[section]

\begin{document}

% === 标题区域（单栏） ===
\twocolumn[
\begin{@twocolumnfalse}
\centering
{\zihao{2}\heiti [论文标题]\par}
\vspace{1em}
{\zihao{4}\songti [作者1]$^{1}$\quad [作者2]$^{2}$\par}
\vspace{0.5em}
{\zihao{5}\songti
$^{1}$[单位1，城市 邮编]\par
$^{2}$[单位2，城市 邮编]\par
}
\vspace{1em}

\begin{minipage}{0.9\textwidth}
\zihao{5}
\textbf{摘\hspace{1em}要\hspace{1em}}[中文摘要，200-300字]

\vspace{0.5em}
\textbf{关键词\hspace{1em}}[关键词1]；[关键词2]；[关键词3]

\vspace{1em}
\textbf{Abstract\hspace{1em}}[English abstract, 150-250 words]

\vspace{0.5em}
\textbf{Keywords\hspace{1em}}keyword1; keyword2; keyword3
\end{minipage}
\vspace{1.5em}
\end{@twocolumnfalse}
]

% === 正文 ===
\input{sections/1_introduction}
\input{sections/2_related_work}
\input{sections/3_method}
\input{sections/4_experiments}
\input{sections/5_discussion}
\input{sections/6_conclusion}

% === 参考文献 ===
\bibliography{references}

\end{document}
